% 本章不引入新的参考文献；在线工具包地址保留为原文脚注。

\section{立体校正算法概要}
\label{sec:rectification-algorithm}

鉴于双目视觉在研究和应用中的广泛使用，我们努力使算法尽可能易于复现和使用。为此，
给出该算法的可运行 MATLAB 代码；代码简洁紧凑（共 22 行），其中的注释使读者无需了解
MATLAB 也能理解代码。\texttt{rectify} 函数（见 MATLAB 代码）的使用方式如下：

\begin{itemize}
  \item 给定一对图像 $I_1,I_2$ 以及通过标定得到的 PPM $P_{o1},P_{o2}$；
  \item 计算 $[T_1,T_2,P_{n1},P_{n2}]=\texttt{rectify}(P_{o1},P_{o2})$；
  \item 应用 $T_1$ 和 $T_2$ 对图像进行立体校正。
\end{itemize}

\begin{lstlisting}[language=Matlab]
function [T1,T2,Pn1,Pn2] = rectify(Po1,Po2)

% RECTIFY: compute rectification matrices

% factorize old PPMs
[A1,R1,t1] = art(Po1);
[A2,R2,t2] = art(Po2);

% optical centers (unchanged)
c1 = - inv(Po1(:,1:3))*Po1(:,4);
c2 = - inv(Po2(:,1:3))*Po2(:,4);

% new x axis (= direction of the baseline)
v1 = (c1-c2);
% new y axes (orthogonal to new x and old z)
v2 = cross(R1(3,:)',v1);
% new z axes (orthogonal to baseline and y)
v3 = cross(v1,v2);

% new extrinsic parameters
R = [v1'/norm(v1)
     v2'/norm(v2)
     v3'/norm(v3)];
% translation is left unchanged

% new intrinsic parameters (arbitrary)
A = (A1 + A2)./2;
A(1,2)=0; % no skew

% new projection matrices
Pn1 = A * [R -R*c1 ];
Pn2 = A * [R -R*c2 ];

% rectifying image transformation
T1 = Pn1(1:3,1:3)* inv(Po1(1:3,1:3));
T2 = Pn2(1:3,1:3)* inv(Po2(1:3,1:3));

% ------------------------
function [A,R,t] = art(P)
% ART: factorize a PPM as  P=A*[R;t]

Q = inv(P(1:3, 1:3));
[U,B] = qr(Q);

R = inv(U);
t = B*P(1:3,4);
A = inv(B);
A = A ./A(3,3);
\end{lstlisting}

一个包含 C 语言和 MATLAB 算法实现、数据集及文档的“立体校正工具包”可在线获取
\footnote{\url{http://www.sci.univr.it/~fusiello/rect.html}}。
